% =====================================================================
% APPENDIX A — Multi-source fused product validation against INMET
% Author: Angela Cunha Soares
% Generated: 2026-05-07
% Sources:
%   output/inmet_validation/per_station_summary.csv
%   output/inmet_validation/pooled_summary.csv
%   output/paper_tables/table_inmet_validation.tex
% =====================================================================

\refstepcounter{section}
\phantomsection
\section*{APPENDIX \thesection: VALIDATION OF THE NASA POWER + OPEN-METEO FUSED CLIMATE PRODUCT AGAINST IN SITU INMET STATIONS (JAN--MAY 2026)}
\label{app:inmet_validation}
\addcontentsline{toc}{section}{\textbf{APPENDIX \thesection: VALIDATION OF THE NASA POWER + OPEN-METEO FUSED CLIMATE PRODUCT AGAINST IN SITU INMET STATIONS}}

\subsection*{A.1.\quad Rationale}

The 1961--2025 backtest of Section~\ref{sec:methods_forecast} is driven
by the Brazilian Daily Weather Gridded Data of \citet{Xavier2022},
which provides the only continuous high-quality precipitation and
ETo series spanning 65 years for the MATOPIBA hubs. Operational
deployment after the publication cutoff requires a substitute
because (i) Xavier is not updated in real time, and (ii) the curated
product cannot be queried from a farm-management software. The
candidate replacement is the **multi-source fusion** of NASA POWER
(MERRA-2, 1990--present) and Open-Meteo Archive (ERA5, 1940--present),
augmented with Open-Meteo Forecast (best-match NWP, today $\pm$ 5 days)
in the operational horizon. Per-variable weights for the fusion follow
the EVAOnline calibration of \citet{Ribeiro2024} against BR-DWGD at
17 sites in Brazil.

This appendix verifies the transfer of those weights to MATOPIBA by
comparing the daily fused product to in situ measurements from the
INMET automatic station network during the only window where direct
ground-truth data are available at the time of writing (1 January
2026 to 7 May 2026, 127 days).

\subsection*{A.2.\quad Stations and coverage}

Five INMET automatic stations sit within or directly adjacent to the
ten MATOPIBA cities used in this study; the conventional station 82768
in Balsas/MA samples only three times daily and is excluded.
Table~\ref{tab:inmet-stations} lists the active stations and their
post-QC coverage for January--May 2026.

\begin{table}[t]
\centering
\caption{INMET automatic stations used in the fused-product validation.
Coverage is the percentage of hourly records with a non-null reading
in the Jan--May 2026 file before any per-variable filtering. Stations
A346, A375 and A416 cleared the per-variable QC threshold of $\geq 30$
days and contribute dedicated columns in
Table~\ref{tab:inmet-validation}; A402 and A404 contribute only via
the across-station pooled row.}
\label{tab:inmet-stations}
\begin{tabular}{llcrrrrr}
\toprule
Code & Municipality & UF & Lat ($^\circ$) & Lon ($^\circ$) & Elev (m)
& Hourly cov. & Days kept \\
\midrule
A346 & Urucu\'i                  & PI & $-7.44$  & $-44.34$ & 399 & 98\,\% & 85 \\
A375 & Baixa Grande do Ribeiro  & PI & $-8.33$  & $-45.09$ & 519 & 97\,\% & 126 \\
A416 & Correntina               & BA & $-13.33$ & $-44.62$ & 552 & 99\,\% & 126 \\
A402 & Barreiras                & BA & $-12.12$ & $-45.03$ & 474 & 25\,\% & 36 (pooled only) \\
A404 & Lu\'is Eduardo Magalh\~aes & BA & $-12.08$ & $-45.71$ & 748 & 52\,\% & 113 (pooled only) \\
\bottomrule
\end{tabular}
\end{table}

\subsection*{A.3.\quad Hourly-to-daily aggregation and quality control}

Hourly INMET records were aggregated to the local civil day (UTC$-3$)
following FAO-56 conventions: $T_\mathrm{max}$ and $T_\mathrm{min}$ as
the per-day extrema of the hourly maximum/minimum readings;
$\overline{T}$, $\overline{\mathrm{RH}}$ and $\overline{u_{10}}$ as
hourly means; and $R_\mathrm{s}$ as the daily sum of incoming
shortwave radiation, converted from kJ\,m$^{-2}$ to MJ\,m$^{-2}$\,d$^{-1}$.
Wind at 10\,m was converted to the 2\,m reference height through
FAO-56 Eq.\,47, $u_2 = u_{10}\cdot 4.87/\ln(67.8\cdot 10 - 5.42)$,
which evaluates to $u_2 \approx 0.748\,u_{10}$ and matches the
``approximately 0.75'' rule of thumb of \citet{Allen1998}.
Reference evapotranspiration was then computed on the cleaned daily
series using the same FAO-56 Penman-Monteith routine
(\verb|compute_eto_fao56_pm|) that is applied downstream of the fusion
operator, so that any structural difference in this comparison
attaches to the climate inputs rather than to the ET formulation.

For each variable a daily aggregate enters the comparison only when
its hourly coverage clears a sensor-specific threshold:
$\geq 20$\,h for $T$, $\geq 18$\,h for RH, $\geq 20$\,h for wind, and
$\geq 8$\,h for $R_\mathrm{s}$ (only daylight hours carry non-null
radiation). Days with fewer valid hours are masked variable-by-variable;
a station enters the per-station column of
Table~\ref{tab:inmet-validation} only if $\geq 4$ variables retain
$\geq 30$ daily aggregates.

\subsection*{A.4.\quad An anomaly in the INMET wind record}

For all three high-coverage stations, $\geq 70\,\%$ of the hourly
anemometer records register exactly $0$\,m\,s$^{-1}$, dragging the
daily mean to roughly $0.15$\,m\,s$^{-1}$. Reference wind speeds for
Cerrado/MATOPIBA at 2\,m sit between 1.5 and 3\,m\,s$^{-1}$
\citep{Allen1998}, so the per-station $u_2$ row in
Table~\ref{tab:inmet-validation} compares the alternative sources
against a degenerate ground truth. We retain that row in the table
to make the issue visible (marked with a dagger), but the
informative wind metric is the pooled row, which compares the
\emph{within-station anomaly} of each daily reading and is therefore
immune to the station-level offset.

\subsection*{A.5.\quad Result --- fused product wins on 5 of 7 variables}

Table~\ref{tab:inmet-validation} (located in
Section~\ref{sec:methods_data}) reports per-station and pooled RMSE for
NASA POWER, Open-Meteo Archive and the EVAOnline-weighted Fused
product against the INMET reference. The fused product attains the
lowest pooled RMSE for $\mathrm{RH}$, $R_\mathrm{s}$, $P$, $u_2$
(anomaly) and $\mathrm{ET}_0$. Open-Meteo Archive (ERA5) wins for
$T_\mathrm{max}$ and $T_\mathrm{min}$ alone, by a small margin
($\Delta$RMSE $< 0.18\,^\circ$C), reflecting that
$\mathrm{HIST}_\mathrm{NASA}$ assigns those two variables a weight
near 0.5; the fusion's improvement is therefore small, and the noise
of a 127-day window can swing it. Crucially, ET$_0$ --- the variable
that drives the irrigation forecast --- is best reproduced by the
fused product (pooled RMSE $0.54$ vs.\ $0.58$ for Archive and
$0.60$ for NASA POWER, a $-7\,\%$ to $-10\,\%$ reduction).

\subsection*{A.6.\quad Comparison against reduced-input ET formulas}

To put the residual disagreement of $0.54$\,mm\,d$^{-1}$ in context,
Table~\ref{tab:eto-alternatives} reports the corresponding daily
RMSE of two reduced-input formulas commonly recommended when the
full Penman-Monteith inputs are unavailable: Hargreaves-Samani
\citep{Hargreaves1985} and Benavides-Lopez \citep{Benavides1970}.
On the same 10 MATOPIBA cities and the longer 1961--2025 Xavier
window, the median RMSEs are $1.08$\,mm\,d$^{-1}$ for
Hargreaves-Samani and $1.37$\,mm\,d$^{-1}$ for Benavides-Lopez ---
two to three times larger than the disagreement between the fused
product and INMET. This indicates that the operational pipeline
keeps the full Penman-Monteith equation, recomputed on the fused
inputs, rather than substituting a reduced formula even though the
fused inputs are themselves estimates rather than direct measurements.

\subsection*{A.7.\quad Limitations}

\begin{itemize}
  \item \textbf{Window is short.} 127 days, January to early May 2026.
  Inter-annual variability and dry-season behaviour are not exercised.
  \item \textbf{Spatial coverage is partial.} Three of ten MATOPIBA
  cities have a dedicated column; two more contribute via the pooled
  row; five (Bom Jesus, Campos Lindos, Formosa do Rio Preto, Tasso
  Fragoso and Balsas) have no INMET automatic station and are not
  represented in this validation.
  \item \textbf{Point-vs-grid representativity.} INMET stations are
  point measurements; NASA POWER reports on a $0.5^\circ\times
  0.625^\circ$ grid; Open-Meteo Archive interpolates to the requested
  coordinate from a $0.25^\circ$ native grid. Some of the residual
  RMSE is representativity error, not source error.
  \item \textbf{The 90-day soybean cycle does not fit inside the
  INMET window.} A typical 1 December planting date predates the
  INMET coverage start of 3 January 2026, so a cycle-level
  verification of the rolling 5-day forecast against INMET ground
  truth is not feasible in the current window. We list it as future
  work (Section~\ref{sec:discussion_limits}) for the
  2026/2027 harvest, which will fall entirely inside the next INMET
  release window.
\end{itemize}

\subsection*{A.8.\quad Reproducibility}

The validation pipeline is implemented in
\verb|scripts/inmet_validation_multistation.py|. Per-station daily
aggregates with QC counters are written to
\verb|output/inmet_validation/per_station/|, the side-by-side aligned
comparison files are in the same folder, and the per-variable
metrics CSVs feeding Table~\ref{tab:inmet-validation} are at
\verb|output/inmet_validation/per_station_summary.csv| and
\verb|output/inmet_validation/pooled_summary.csv|.
The raw INMET hourly downloads are in \verb|data_raw/inmet/|.
